\section{$F$-гомоморфизмы простых расширений}

\begin{stmt}
    Пусть $F(\delta)$ --- простое расширение, $\varphi\colon F \to L$ --- гомоморфизм полей, и $\tilde{\varphi}
    \colon F(\delta) \to L$ --- его продолжение (то есть $\tilde{\varphi}\big|_F = \varphi$). Тогда для любого
    $p(x) \in F[x]$:
    \[
        \tilde{\varphi}(p(\delta)) = p(\tilde{\varphi}(\delta)),
    \]
    где в правой части коэффициенты $p$ заменяются на их образы при $\varphi$.

    \begin{proof}
        Пусть $p(x) = d_n x^n + d_{n-1} x^{n-1} + \dots + d_0$. Тогда:
        \begin{align*}
            \tilde{\varphi}(p(\delta)) &= \tilde{\varphi}(d_n \delta^n + d_{n-1}\delta^{n-1} + \dots + d_0)
            \\
            &= \tilde{\varphi}(d_n)\,\tilde{\varphi}(\delta)^n + \tilde{\varphi}(d_{n-1})\,\tilde{\varphi}
            (\delta)^{n-1} + \dots + \tilde{\varphi}(d_0) \\
            &= d_n\,\tilde{\varphi}(\delta)^n + d_{n-1}\,\tilde{\varphi}(\delta)^{n-1} + \dots + d_0 = p(\tilde{\varphi}
            (\delta)),
        \end{align*}
        так как $\tilde{\varphi}(d_i) = \varphi(d_i) = d_i$ при $\varphi\big|_F = \mathrm{id}$.
    \end{proof}
\end{stmt}


\begin{example}
    $F = \mathbb{R}$, $\delta = i$, $F(\delta) = \mathbb{R}(i) = \mathbb{C}$. Рассмотрим $\mathbb{R}$-гомоморфизмы
    $\varphi\colon \mathbb{C} \to \mathbb{C}$.

    Так как $p(x) = x^2 + 1$ --- минимальный многочлен $i$, то $\varphi(i)$ --- корень $x^2 + 1$, то есть
    $\varphi(i) = \pm i$.

    Таких гомоморфизмов ровно $2$: тождественный ($\varphi(i) = i$) и комплексное сопряжение ($\varphi(i)
    = -i$).
\end{example}


\section{Алгебраические элементы и простые расширения}

\begin{definition}
    Пусть $L/F$ --- расширение полей. Элемент $d \in L$ называется \emph{алгебраическим} над $F$, если
    $\exists\, p(x) \in F[x]\colon p(d) = 0$.
\end{definition}


\begin{stmt}
    Пусть $d \in L$ --- алгебраический над $F$. Тогда существует единственный унитарный неприводимый многочлен
    $p(x) \in F[x]$ минимальной степени, аннулирующий $d$ --- \emph{минимальный многочлен} $d$ над $F
    $.

    \begin{proof}
        Пусть $p(x)$ --- минимальный. Если $p = q_1 \cdot q_2$, то $0 = p(d) = q_1(d) \cdot q_2(d)$, откуда
        $q_1(d) = 0$ или $q_2(d) = 0$ (так как $L$ --- поле). Противоречие с минимальностью степени. Значит
        $p$ неприводим.

        Следовательно, $(p(x))$ --- максимальный идеал в $F[x]$, и
        \[
            F[d] \cong \frac{F[x]}{(p(x))} \text{ --- поле.}
        \]
    \end{proof}
\end{stmt}


\begin{stmt}
    Если $d$ алгебраичен над $F$ и $n = \deg p(x)$, то $F[d]$ --- конечномерное пространство над $F$ с базисом
    $1, d, d^2, \dots, d^{n-1}$. В частности, $[F[d] : F] = \deg p(x)$.
\end{stmt}


\begin{examplelist}
    \item $\mathbb{R}[i] = \mathbb{R}[x]/(x^2 + 1) = \mathbb{C}$, базис над $\mathbb{R}$: $\{1, i\}$.


    \item $\mathbb{Q}[\sqrt[5]{2}]$ --- расширение степени $5$ с базисом $\{1, \sqrt[5]{2}, (\sqrt[5]
    {2})^2, (\sqrt[5]{2})^3, (\sqrt[5]{2})^4\}$.

    \item $\delta = \sqrt{2} + \sqrt{3}$. Вычислим: $\delta^2 = 5 + 2\sqrt{6}$, откуда $(\delta^2 - 5)
    ^2 = 24$, то есть $\delta$ --- корень $(x^2 - 5)^2 - 24 = 0$. Значит $\delta$ алгебраичен над $\mathbb{Q}
    $.
\end{examplelist}


\begin{stmt}
    Если $d$ алгебраичен над $F$, то $F[d] = F(d)$, то есть кольцо $F[d]$ уже является полем.
\end{stmt}

\begin{remark}
    Если $d$ трансцендентен, то $F[d] \cong F[x]$ --- кольцо многочленов, а не поле. Поле частных:
    \[
        F(d) = \left\{\frac{p(d)}{q(d)} \;\middle|\; p, q \in F[x],\, q(d) \neq 0\right\}.
    \]
    Например, $\mathbb{Q}(\pi) \cong \mathbb{Q}(x)$ --- поле рациональных функций.
\end{remark}


\begin{stmt}
    $F[d, \beta] = F[d][\beta]$ --- пространство конечной размерности над $F$, если $d$ и $\beta$ алгебраичны.

\end{stmt}


\section{Мультипликативность степеней расширений}

\begin{stmt}
    Пусть $F \subset T \subset L$ --- башня расширений. Тогда
    \[
        [L : F] = [L : T] \cdot [T : F].
    \]
\end{stmt}


\section{Гомоморфизмы алгебраических расширений}

\begin{stmt}
    Гомоморфизм полей $\varphi\colon F \to L$ всегда инъективен.

    \begin{proof}
        $\ker\varphi$ --- идеал в $F$. Если $\ker\varphi \neq 0$, то $\exists\, x \in \ker\varphi$, $x \neq 0
        $.
        Тогда $1 = x^{-1} \cdot x \in \ker\varphi$, откуда $\ker\varphi = F$ и $\varphi(F) = 0$. Но $\varphi(1)
        = 1 \neq 0$ --- противоречие.
        Значит $\ker\varphi = 0$.
    \end{proof}
\end{stmt}


\begin{stmt}
    Пусть $f \in F[x]$, $d_1, \dots, d_n$ --- все корни $f$, $E = F[d_1, \dots, d_n]$ --- поле разложения
    $f$.
    Пусть $L/F$ --- расширение, содержащее поле разложения $f$.

    Тогда число $F$-гомоморфизмов $\varphi\colon E \to L$ не превосходит $[E : F]$. Оно равно $[E : F]
    $, если все корни $f$ различны.

    \begin{proof}
        Пусть $E = F[d_1, \dots, d_n]$. Для каждого $d_i$ возьмём минимальный многочлен $f_i(x)$ над
        $F$. Тогда $f_i \mid f$ и $f_i$ неприводим.

        Строим $\varphi$ индуктивно:

        \textbf{Шаг 1.} Число $F$-гомоморфизмов $\varphi_1\colon F[d_1] \to L$ равно числу корней $f_1(x)
        $ в $L$, что не превосходит $\deg f_1 = [F[d_1] : F]$.

        \textbf{Шаг 2.} Для каждого $\varphi_1$ ищем продолжение $\varphi_2\colon F[d_1, d_2] \to L$.

        Рассматриваем минимальный многочлен $d_2$ над $F[d_1]$; он делит $f(x)$ в $F[d_1][x]$ и имеет различные корни,

        если $f(x)$ имеет различные корни. Таких продолжений $\leq [F[d_1, d_2] : F[d_1]]$.

        \textbf{Продолжая}, получаем $\varphi_n\colon E \to L$, и число таких гомоморфизмов:
        \[
            \leq [F[d_1] : F] \cdot [F[d_1, d_2] : F[d_1]] \cdot \ldots = [E : F].
        \]
        Равенство достигается, когда на каждом шаге все корни различны.
    \end{proof}
\end{stmt}


\begin{cor}
    Пусть $E_1$, $E_2$ --- два минимальных поля разложения $f(x) \in F[x]$. Тогда любой $F$-гомоморфизм
    $E_1 \to E_2$ --- изоморфизм.

    \begin{proof}
        $F$-гомоморфизм $\varphi\colon E_1 \to E_2$ инъективен (как гомоморфизм полей), поэтому $[E_1 :
        F] \leq [E_2 : F]$.

        По предложению выше, $\exists\, \psi\colon E_2 \to E_1$, и аналогично $[E_2 : F] \leq [E_1 : F]
        $.

        Значит $[E_1 : F] = [E_2 : F]$, и $\varphi$, $\psi$ --- изоморфизмы.
    \end{proof}
\end{cor}


\begin{cor}
    Пусть $E$ и $L$ --- расширения $F$, $E/F$ конечное. Тогда число $F$-гомоморфизмов $E \to L$ не превосходит
    $[E : F]$.
\end{cor}


\section{Нормальные расширения}

\begin{definition}
    Алгебраическое расширение $E/F$ называется \emph{нормальным}, если $\forall\, d \in E$ минимальный многочлен
    $p(x)$ элемента $d$ над $F$ раскладывается на линейные множители в $E[x]$.

    Эквивалентно: $\forall$ неприводимый $f(x) \in F[x]$, имеющий хотя бы один корень в $E$, имеет в
    $E$ все свои корни.
\end{definition}


\begin{example}[Ненормальное расширение]
    $\mathbb{Q}(\sqrt[3]{2}) / \mathbb{Q}$ --- ненормальное. Минимальный многочлен $\sqrt[3]{2}$ --- это
    $x^3 - 2$, который имеет три корня: $\sqrt[3]{2}$, $\sqrt[3]{2}\,\omega$, $\sqrt[3]{2}\,\omega^2$,
    где $\omega = e^{2\pi i/3}$. Но $\sqrt[3]{2}\,\omega \notin \mathbb{Q}(\sqrt[3]{2}) \subset \mathbb{R}
    $.
\end{example}


\section{Сепарабельные расширения}

\begin{definition}
    Многочлен $p(x) \in F[x]$ называется \emph{сепарабельным}, если все его корни (в поле разложения)
    простые.

    Элемент $d$ расширения $E/F$ называется \emph{сепарабельным}, если его минимальный многочлен над
    $F$ сепарабелен.

    Алгебраическое расширение $E/F$ называется \emph{сепарабельным}, если каждый $d \in E$ сепарабелен.

\end{definition}


\begin{example}[Несепарабельное расширение]
    Пусть $k = \mathbb{F}_p$. Рассмотрим $\mathbb{F}_p(\delta)$ над $\mathbb{F}_p(\delta^p)$. Минимальный многочлен
    $\delta$ над $\mathbb{F}_p(\delta^p)$ --- это $x^p - \delta^p$, и его производная $(x^p - \delta^p)
    ' = p x^{p-1} = 0$. Значит у многочлена кратные корни, и расширение несепарабельно.
\end{example}


\section{Расширения Галуа}

\begin{definition}
    Конечное алгебраическое расширение $E/F$ называется \emph{расширением Галуа}, если оно одновременно сепарабельно и нормально.

\end{definition}

\begin{remark}
    Нормальность и сепарабельность означают: если $E/F$ --- расширение Галуа и $f(x) \in F[x]$ --- неприводимый многочлен,
    имеющий хотя бы один корень в $E$, то все его корни лежат в $E$, и их ровно $\deg f$ (все различны).
\end{remark}
